\section{Đánh giá, phân công và kết luận}

\subsection{Đánh giá phương pháp}

Phương pháp CAFE (Cognition of humAn for Facial Expression) đã thể hiện những ưu điểm vượt trội về mặt lý thuyết lẫn thực nghiệm trong bài toán nhận dạng cảm xúc khuôn mặt có khả năng tổng quát hóa cao (Generalizable FER). Qua quá trình nghiên cứu và thực nghiệm tái hiện kết quả, nhóm đưa ra các đánh giá kỹ thuật trọng tâm sau:

\begin{itemize}
    \item \textbf{Cơ chế phân tách đặc trưng thông minh:} Bằng cách đóng băng toàn bộ trọng số của mô hình ngôn ngữ - hình ảnh lớn CLIP ($ViT-B/32$), phương pháp đã tận dụng được không gian biểu diễn khuôn mặt cực kỳ tổng quát được học từ hàng triệu ảnh Internet. Việc này giúp mô hình hoàn toàn miễn nhiễm với hiện tượng học quá khớp (overfitting) các đặc trưng nhiễu cục bộ của tập dữ liệu nguồn như điều kiện ánh sáng, pose mặt hay đặc trưng danh tính cá nhân.
    \item \textbf{Mặt nạ Sigmoid thích ứng hiệu quả:} Việc sử dụng mạng ResNet-18 gọn nhẹ làm bộ tạo mặt nạ $M_s = \text{Sigmoid}(M)$ cho phép mô hình học cách tinh lọc động các kênh đặc trưng liên quan đến cảm xúc cho từng ảnh đầu vào. Cơ chế này hoạt động như một bộ lọc xác suất mềm, giữ nguyên các thông tin hữu ích và triệt tiêu các kênh dư thừa của CLIP mà không làm biến dạng cấu trúc đặc trưng gốc.
    \item \textbf{Thiết kế bộ phân loại không phụ thuộc lớp liên kết đầy đủ (FC):} Đây là một cải tiến táo bạo nhằm loại bỏ hoàn toàn tầng FC ở phần cuối mạng --- vốn là tác nhân chính gây ra việc ghi nhớ nhãn cục bộ. Việc chia cắt trực tiếp khối đặc trưng đã lọc $\tilde{F}$ thành $L=7$ phần dọc theo chiều kênh và áp dụng phép toán gộp cực đại (max pooling) giúp mô hình tối ưu hóa trực tiếp trên ranh giới quyết định của từng lớp cảm xúc.
    \item \textbf{Độ bền bỉ và tính hội tụ cao:} Hàm mất mát đa nhiệm kết hợp tuyến tính $l_{train} = l_{cls} + \lambda l_{sep} + \beta l_{div}$ giúp ràng buộc chặt chẽ tính phân tách kênh ($l_{sep}$) và tính đa dạng kênh ($l_{div}$). Kết quả thực nghiệm của nhóm trên tập RAF-DB đạt độ chính xác $89.15\%$ (vượt $0.43\%$ so với bài báo gốc) và kết quả đánh giá chéo trên AffectNet đạt $52.01\%$ (vượt $6.15\%$ so với tham chiếu) đã chứng minh độ bền bỉ vượt trội của phương pháp ngay cả khi thay đổi cấu trúc mạng xương sống.
\end{itemize}

\subsection{Hạn chế và vấn đề chưa giải quyết}

Mặc dù đạt được những kết quả khả quan về khả năng tổng quát hóa zero-shot, thực nghiệm chi tiết của nhóm trên các tập dữ liệu ngoài miền (out-of-domain) đã chỉ ra một số hạn chế cốt lõi sau:

\begin{itemize}
    \item \textbf{Sự sụt giảm nghiêm trọng hiệu suất trên các lớp ít mẫu (Class Imbalance):} Phân tích ma trận nhầm lẫn trên các tập dữ liệu như CK+48, FERPlus, và SFEW2.0 cho thấy mô hình bị thiên lệch nặng nề về các lớp chiếm đa số như \textit{happy} và \textit{neutral}. Ngược lại, các lớp cảm xúc có tần suất xuất hiện thấp và có biểu hiện cơ mặt phức tạp như \textit{fear} (sợ hãi), \textit{disgust} (ghê tởm) có độ chính xác cực kỳ thấp, thậm chí đạt $0\%$ trên tập CK+48 và SFEW2.0. Hàm mất mát hiện tại chưa có cơ chế trọng số để bù trừ sự mất cân bằng này.
    \item \textbf{Nhạy cảm với chất lượng và độ phân giải hình ảnh:} Đối với tập dữ liệu CK+48, do ảnh gốc có độ phân giải rất thấp ($48 \times 48$ pixel) nên khi phóng đại lên $224 \times 224$ để phù hợp đầu vào mô hình, các đặc trưng hình học và biểu cảm cơ mặt vi tế đã bị mờ nhòe. Do CLIP được huấn luyện trên ảnh chất lượng cao, sự không khớp về chất lượng ảnh khiến mặt nạ Sigmoid khó định vị chính xác vùng đặc trưng cảm xúc.
    \item \textbf{Giới hạn của thông tin ảnh tĩnh (Static FER):} Mô hình CAFE chỉ xử lý thông tin trên từng khung ảnh độc lập. Trong thực tế, các biểu cảm khuôn mặt mang tính động và diễn tiến theo thời gian. Việc thiếu đi ngữ cảnh thời gian khiến mô hình dễ nhầm lẫn giữa các trạng thái cảm xúc chuyển tiếp hoặc mơ hồ về mặt ngữ nghĩa (ví dụ giữa \textit{sadness} và \textit{neutral}, hoặc \textit{fear} và \textit{surprise}).
    \item \textbf{Sự phụ thuộc vào phân phối nhãn nguồn:} Mô hình được huấn luyện trên nhãn cứng của RAF-DB vốn được gán bởi các chuyên gia dựa trên phân phối chuẩn. Khi chuyển sang các tập dữ liệu có quy ước gán nhãn khác (như FERPlus sử dụng cơ chế bỏ phiếu đám đông để tạo nhãn phân phối), sự sai lệch trong định nghĩa biên cảm xúc làm giảm độ chính xác tổng thể.
\end{itemize}

\subsection{Nhận định cá nhân và hướng phát triển}

Từ góc nhìn của những người thực hiện đồ án, nhóm nhận thấy phương pháp CAFE đã mở ra một tư duy thiết kế rất thực tiễn cho các ứng dụng nhận dạng cảm xúc trong thế giới thực. Việc loại bỏ hoàn toàn yêu cầu thu thập dữ liệu miền đích giúp mô hình dễ dàng tích hợp trực tiếp vào các API thương mại hoặc các hệ thống nhúng (như camera giám sát người lái xe) mà không lo ngại về vấn đề quyền riêng tư dữ liệu hay chi phí tinh chỉnh lại mạng.

Để khắc phục các hạn chế hiện hữu và nâng cao hơn nữa hiệu suất của hệ thống, nhóm đề xuất một số hướng phát triển và cải tiến kỹ thuật tiềm năng trong tương lai:

\begin{enumerate}
    \item \textbf{Tích hợp hàm mất mát cân bằng lớp (Class-Balanced Loss):} Thay thế hàm Cross-Entropy thông thường bằng Focal Loss hoặc Cost-Sensitive Loss kết hợp với bộ lấy mẫu cân bằng (Balanced Sampler) trong quá trình huấn luyện trên RAF-DB. Điều này sẽ ép buộc mô hình tập trung học các đặc trưng của các lớp thiểu số như \textit{fear} và \textit{disgust}, nâng cao chỉ số Mean Accuracy và Macro-F1.
    \item \textbf{Đào tạo đa độ phân giải thích ứng (Multi-resolution Training):} Tăng cường dữ liệu huấn luyện bằng cách làm mờ (blur) hoặc hạ độ phân giải ngẫu nhiên của tập RAF-DB để mô hình học được các đặc trưng bất biến với chất lượng hình ảnh, từ đó cải thiện hiệu suất kiểm thử trên các tập ảnh chất lượng thấp như CK+48.
    \item \textbf{Mở rộng sang nhận dạng biểu cảm động (Video-based FER):} Phát triển mô-đun temporal adapter dựa trên cấu trúc mạng hồi quy GRU, LSTM hoặc mạng tự chú ý Transformer đặt phía sau khối đặc trưng $\tilde{F}$. Sự bổ sung này sẽ giúp mô hình bắt giữ được sự thay đổi cơ mặt theo thời gian, giải quyết triệt để sự mơ hồ giữa các lớp cảm xúc tĩnh.
    \item \textbf{Hướng dẫn mặt nạ bằng gợi ý văn bản cấu trúc (Text-guided Masking):} Tận dụng bộ mã hóa văn bản (Text Encoder) của CLIP để đưa vào các mô tả chi tiết về Đơn vị hành động cơ mặt (Action Units --- FACS) tương ứng với từng loại cảm xúc. Việc kết hợp thông tin đa phương thức (multi-modal prompting) này hứa hẹn sẽ định hướng mặt nạ Sigmoid hội tụ vào các vùng cơ mặt đích một cách chính xác và tường minh hơn.
\end{enumerate}

\subsection{Kế hoạch và phân công nhiệm vụ}
\begin{table}[H]
    \centering
    \caption{Phân công nhiệm vụ và tỷ lệ hoàn thành}
    \begin{tabularx}{\textwidth}{C{2.3cm} L{3.2cm} X C{2.0cm}}
        \toprule
        \textbf{MSSV} & \textbf{Thành viên} & \textbf{Nhiệm vụ phụ trách} & \textbf{Tỷ lệ hoàn thành} \\
        \midrule

        22120134 & Hoàng Tiến Huy 
        & Tìm hiểu mã nguồn của phương pháp SOTA, phân tích cơ sở thiết kế của hệ thống và triển khai demo nhận diện cảm xúc theo thời gian thực bằng webcam.
        & 100\% \\

        23120195 & Lê Hà Thanh Chương
        & Khảo sát một bài báo SOTA về nhận diện cảm xúc, phân tích động lực nghiên cứu, ý tưởng cốt lõi, điểm khác biệt so với giáo trình và đưa ra nhận xét cá nhân.
        & 100\% \\

        23120232 & Lê Thượng Đế
        & Nghiên cứu nội dung Facial Expression trong giáo trình, hệ thống hóa nền tảng lý thuyết, quy trình xử lý và các công thức toán học liên quan.
        & 100\% \\

        23120245 & Nguyễn Quang Duy
        & Tổng hợp nội dung và kết quả của nhóm, rà soát tính thống nhất của báo cáo, xây dựng slide trình bày và chuẩn bị kịch bản thuyết trình.
        & 100\% \\

        23120258 & Lưu Trọng Hiếu
        & Ghi nhận kết quả thực nghiệm từ mã nguồn, đối chiếu với kết quả trong bài báo, phân tích nguyên nhân sai lệch và hoàn thiện phần nhận xét thực nghiệm.
        & 100\% \\

        23120260 & Văn Đình Hiếu
        & Cài đặt môi trường thực nghiệm, chạy lại mã nguồn của bài báo và thực hiện thêm thử nghiệm trên dataset ngoài để đánh giá khả năng tổng quát của mô hình.
        & 100\% \\

        \bottomrule
    \end{tabularx}
\end{table}

\subsection{Kết luận}

Đồ án môn học ``Nhận dạng mẫu'' với chủ đề Nhận dạng cảm xúc khuôn mặt có khả năng tổng quát hóa cao (Generalizable FER) đã hoàn thành trọn vẹn tất cả các mục tiêu đề ra ban đầu trên cả hai phương diện nghiên cứu lý thuyết và thực hành thực nghiệm. Thông qua quá trình triển khai dự án, nhóm đã đạt được các kết quả cụ thể sau:

\begin{itemize}
    \item \textbf{Về mặt lý thuyết:} Nhóm đã hệ thống hóa thành công các kiến thức nền tảng về quy trình của một hệ thống phân tích biểu cảm khuôn mặt tự động (AFEA) truyền thống từ giáo trình chuyên ngành, đồng thời nghiên cứu chuyên sâu bài báo khoa học SOTA tại hội nghị ECCV 2024 về phương pháp CAFE. Sự đối sánh kỹ thuật chi tiết giữa các đặc trưng thủ công (geometric, appearance features) và đặc trưng học sâu đa phương thức hiện đại đã giúp nhóm làm rõ bước chuyển dịch công nghệ của lĩnh vực nhận dạng mẫu trong kỷ nguyên mô hình nền tảng.
    \item \textbf{Về mặt thực nghiệm:} Nhóm đã tái cài đặt và huấn luyện thành công mô hình CAFE trên tập dữ liệu nguồn RAF-DB, đồng thời thực hiện đánh giá chéo chéo miền (zero-shot cross-domain) trên năm bộ dữ liệu đích bao gồm FERPlus, AffectNet, SFEW2.0, MMAFEDB và bộ dữ liệu mở rộng ngoài bài báo CK+48. Các kết quả thực nghiệm thu được có độ chính xác rất cao, tiệm cận và thậm chí vượt trội hơn so với số liệu được công bố trong bài báo gốc (đạt độ chính xác tổng thể $89.15\%$ trên RAF-DB và $52.01\%$ trên AffectNet), qua đó xác thực mạnh mẽ tính đúng đắn của giả thuyết nghiên cứu.
    \item \textbf{Về mặt ứng dụng:} Nhóm đã phát triển thành công ứng dụng demo nhận dạng cảm xúc thời gian thực qua webcam laptop theo hai hướng tiếp cận song song: sử dụng mô hình CAFE đã huấn luyện và sử dụng mô hình CLIP thuần zero-shot với các prompt văn bản tùy biến. Cả hai ứng dụng đều vận hành ổn định, trực quan hóa sinh động kết quả dự đoán và điểm số tin cậy trực tiếp trên luồng video của người dùng.
\end{itemize}

Nhìn chung, đồ án này không chỉ giúp các thành viên trong nhóm củng cố vững chắc các kiến thức lý thuyết đã học trên lớp, mà còn rèn luyện kỹ năng giải quyết các bài toán thực tế của học sâu như sự lệch miền dữ liệu (domain gap), chất lượng ảnh thấp và sự mất cân bằng phân phối lớp nghiêm trọng. Những kết quả đạt được và các hướng phát triển được đề xuất trong báo cáo sẽ là nền tảng quan trọng giúp nhóm tiếp tục nghiên cứu sâu hơn về thị giác máy tính và trí tuệ nhân tạo trong tương lai.